% Notes of the Montreal seminar, 1970 (cote 129, typed part): archive cote 129, site page 40 (PDF page 41).
% typescript. Draft compared: none.
% AI-transcribed by Claude Opus 5.5 (claude-opus-5-5, adaptive thinking, effort high), Message Batches API, batch mtl129_full (msgbatch_01AGnDKHwy6ATCM7ehUa6263), 2026-11-13; not yet checked by eye.
\scanpage{129/40}
\mspage{\hand{II.12}}

On a les équivalences :

$G$ est infinitésimal $\Longleftrightarrow G^{\mathrm{et}} = 0$

$G$ est étale $\Longleftrightarrow G^{\mathrm{mult}} = G^{\mathrm{bi}} = 0$

$G$ est unipotent $\Longleftrightarrow G^{\mathrm{mult}} = 0$

$G$ est bi-infinitésimal $\Longleftrightarrow G^{\mathrm{mult}} = G^{\mathrm{et}} = 0$

$G$ est de type multiplicatif $\Longleftrightarrow G^{\mathrm{bi}} = G^{\mathrm{et}} = 0$ .

3.3. La catégorie des p-groupes étales finis sur un corps parfait
$k$ est bien connue. Elle est équivalente à la catégorie des p-groupes fi-%
nis sur lesquels opère le groupe fondamental $\pi_1(k) = \mathrm{Gal}(\overline{k}/k)$ . Par dua-%
lité de Cartier, la catégorie des p-groupes finis de type multiplicatif
est équivalente à la catégorie opposée de celle des p-groupes étales. Par
ailleurs, la catégorie des p-groupes finis sur lesquels opère $\pi_1(k)$ est
équivalente à sa propre opposée par dualité de Prontryagin. On en déduit
une équivalence de catégories entre la catégorie des p-groupes finis de
type multiplicatif avec celle des p-groupes finis sur lesquels opère $\pi_1(k)$ .

4. \uline{Théorie de Dieudonné sur un corps parfait}

4.1. Soit $k$ un corps parfait de caractéristique $p > 0$ et soit
p-Gr-fin(k) la catégorie des schémas en groupes finis sur $k$ et annulés
par une puissance de $p$ . Le but de la théorie de Dieudonné est d'établir
une anti-équivalence entre la catégorie p-Gr-fin(k) et une catégorie de
modules que nous allons définir.

On appelle \uline{anneau de Dieudonné} \hand{$\mathcal{D}$} l'anneau $W[F,V]$ (non com-%
mutatif) engendré sur l'anneau local $W = W(k)$ (I,4) par des indéterminées

\margin{40}
